% theme: acids_alkalis_and_ions
\MajorQuestion{酸・アルカリとイオン}
\SetRowSpaceQanda

\Instruction{次の問いに答えなさい。}
\QQ{水に溶けたとき、電離して水素イオンを生じる物質を\NamedBlank{acid}という。}{酸}
\QQ{BTB溶液を加えると青色になる水溶液は、何性か。}{アルカリ性}
\QQ{酸とアルカリが、たがいの性質を打ち消し合う化学変化を\NamedBlank{neutralization}という。}{中和}
\QQ{水酸化物イオンを、イオンの記号で表しなさい。}{$\mathrm{OH^-}$}
\QQ{pHが7の水溶液は、何性か。}{中性}

\Instruction{\strut}
\QQ{水に溶けたとき、電離して水酸化物イオンを生じる物質を\NamedBlank{alkali}という。}{アルカリ}
\QQ{\RefBlank{neutralization}では、水素イオンと水酸化物イオンが結びついて何ができるか。}{水}
\QQ{青色リトマス紙を赤色に変える水溶液は、何性か。}{酸性}
\QQ{水溶液に電圧を加えると、水素イオンは陰極と陽極のどちらに引かれるか。}{陰極}
\QQ{酸性の水溶液にマグネシウムを入れると発生する気体は何か。}{水素}

\Instruction{\strut}
\QQ{\RefBlank{neutralization}のとき、水のほかにできる物質を\NamedBlank{salt}という。}{塩}
\QQ{pHの値が7より小さいほど、酸性は強くなるか弱くなるか。}{強くなる}
\QQ{塩化水素が水に溶けて電離するようすを、化学式で表しなさい。}{$\mathrm{HCl \rightarrow H^+ + Cl^-}$}
\QQ{赤色リトマス紙を青色に変える水溶液は、何性か。}{アルカリ性}
\QQ{水素イオンを、イオンの記号で表しなさい。}{$\mathrm{H^+}$}

\Instruction{\strut}
\QQ{\RefBlank{salt}は、酸の陰イオンとアルカリの陽イオンが結びついてできる。塩酸と水酸化ナトリウム水溶液からできる\RefBlank{salt}は何か。}{塩化ナトリウム}
\QQ{pHの値が7より大きいほど、アルカリ性は強くなるか弱くなるか。}{強くなる}
\QQ{水溶液に電圧を加えると、水酸化物イオンは陰極と陽極のどちらに引かれるか。}{陽極}
\QQ{BTB溶液を加えると黄色になる水溶液は、何性か。}{酸性}
\QQ{塩酸と水酸化ナトリウム水溶液の中和を、化学反応式で表しなさい。}{$\mathrm{HCl + NaOH \rightarrow NaCl + H_2O}$}
